\chapter*{Kurzfassung}
\addcontentsline{toc}{chapter}{Kurzfassung}

In vielen Automatisierungsanwendungen liegen Bauteile nicht an einer festen Position, sondern
werden mit wechselnder Lage bereitgestellt. Gleichzeitig soll ihre Qualität möglichst direkt im
Prozess geprüft werden. In dieser Projektarbeit wird dazu ein Demonstrator aufgebaut, der beide
Aufgaben verbindet: Ein kollaborativer Roboter vom Typ LARA der Firma Neura Robotics greift ein
Bauteil, das an beliebiger Position und in beliebiger Orientierung in einem Arbeitsbereich liegt,
führt es einer automatischen Sichtprüfung zu und sortiert es anschließend als Gut- oder
Schlechtteil aus.

Das System besteht aus mehreren Modulen, die über fest definierte Schnittstellen miteinander
verbunden sind. Die Objekterkennung findet den Arbeitsbereich im Bild einer Kamera über vier
ArUco-Marker, erzeugt daraus eine entzerrte Draufsicht und bestimmt Position und Orientierung
des Bauteils. Eine abgeschrägte Ecke des Bauteils macht die Orientierung dabei eindeutig. Die
Robotersteuerung rechnet diese Lage in eine Greifpose um und führt die Greif- und Ablegebewegungen
mit reduzierter Geschwindigkeit innerhalb fester Arbeitsraumgrenzen aus. An der Prüfstation
liest die Sichtprüfung mit einer seitlichen Kamera eine fünfstellige Seriennummer mittels
Texterkennung, misst den Durchmesser eines Durchgangslochs über eine Ellipsenanpassung und prüft,
ob eine Kerbe vorhanden ist. Als Alternative zum programmierten Greifen wird ein lernbasierter
Ansatz entwickelt: Eine Diffusion Policy lernt aus Demonstrationen, die der Roboter mit
gezielt eingespeistem Rauschen selbst fährt, das Bauteil anhand der Bilder einer Kamera am
Greifer und der Kamera über dem Arbeitsbereich zu greifen. Sie ist über eine gemeinsame
Schnittstelle in den Ablauf eingebunden und übergibt das Bauteil über der Prüfstation an die
Robotersteuerung.

Ein zentraler Zustandsautomat steuert den Ablauf aus Erkennen, Greifen, Prüfen und Sortieren,
reagiert auf Fehler mit einem Stopp des Roboters und protokolliert jeden Zyklus. Bedient wird
das System über eine grafische Oberfläche. Jedes Modul kann mit Platzhaltern, gespeicherten
Testbildern, einer Nachbildung der Robotersteuerung, der Simulation des Herstellers oder der
realen Hardware betrieben werden,
sodass sich die Module unabhängig voneinander entwickeln und schrittweise in Betrieb nehmen
lassen. Die eingesetzten Verfahren der Bildverarbeitung sind bewusst einfach und
nachvollziehbar gehalten.

Am realen Roboter lief ein erster vollständiger Ablauf durch, mit Erkennung und Prüfung auf
Testbildern. Die Erprobung an der realen Anlage umfasst
wiederholte Durchläufe mit einem Gutteil in zufälliger Lage sowie gezielt fehlerhafte Bauteile
ohne Kerbe, mit abweichendem Lochdurchmesser und ohne Seriennummer.
\todo[inline]{Ergebnisse der Versuche in ein bis zwei Sätzen ergänzen.}
